% TA-PFC-E1 — Luis Yépez. Visión y objetivos de alto nivel.
% Fragmento para insertar después de \subsection{Visión general y mejora esperada con el sistema}.
% La versión del alcance proviene de la actualización comunicada por el equipo.

La visión del proyecto es desarrollar una \textbf{aplicación web para la gestión y seguimiento de tutorías académicas mediante el análisis de indicadores de rendimiento estudiantil, orientada a la detección temprana de dificultades académicas}. Su propósito es ofrecer a estudiantes, docentes tutores y coordinación académica un medio organizado para programar tutorías, registrar el acompañamiento realizado y consultar información pertinente sobre la evolución académica de cada estudiante, de acuerdo con permisos institucionales que deberán definirse.

La plataforma propuesta relacionará el historial de tutorías y asistencias con indicadores de rendimiento que el equipo deberá acordar con los interesados. A partir de esa información, se pretende facilitar la \emph{identificación oportuna de posibles dificultades} y apoyar la recomendación de tutorías para su revisión por responsables académicos. Las alertas o recomendaciones serán apoyos para la intervención humana: no equivalen a un diagnóstico concluyente ni sustituyen las decisiones de docentes o autoridades.

\paragraph{Objetivos académicos y organizacionales de alto nivel.}
\begin{enumerate}
  \item \textbf{OBJ-01. Organizar la atención tutorial:} centralizar solicitudes, programación y consulta de tutorías académicas para estudiantes y docentes.
  \item \textbf{OBJ-02. Registrar el acompañamiento:} conservar las sesiones realizadas y su asistencia, con identificación de la tutoría y de los participantes autorizados.
  \item \textbf{OBJ-03. Visualizar la evolución académica:} presentar indicadores de rendimiento y tendencias de calificaciones disponibles por vías autorizadas, sin alterar los registros oficiales.
  \item \textbf{OBJ-04. Detectar señales tempranas:} identificar posibles dificultades académicas mediante criterios que deberán definirse y validarse con los responsables del proceso tutorial.
  \item \textbf{OBJ-05. Apoyar intervenciones oportunas:} proponer tutorías o acciones de seguimiento para que el personal académico correspondiente las revise y decida su pertinencia.
  \item \textbf{OBJ-06. Resguardar la información:} diferenciar el acceso de estudiantes, docentes, administradores y coordinadores, evitando consultas de datos académicos fuera de los permisos autorizados.
\end{enumerate}

\paragraph{Alcance y supuestos pendientes de validación.}
Esta visión constituye una propuesta del primer avance, no un conjunto de requisitos aprobado por la institución. Deben confirmarse las fuentes de calificaciones, la autorización para consultarlas, los indicadores que serán útiles, los criterios para señalar riesgo académico y quién podrá revisar las recomendaciones. El sistema se plantea como una herramienta de consulta y seguimiento; \textbf{no se propone modificar calificaciones oficiales ni resolver problemas de integridad transaccional de su registro}. La publicación de modelos o materiales de clases no forma parte de las funcionalidades priorizadas en la actualización recibida y solo se incorporaría si el equipo confirma expresamente su inclusión.
